\documentclass[11pt]{article}
\usepackage[a4paper,margin=25mm]{geometry}
\usepackage{amsmath,amssymb,amsthm,hyperref,longtable}
\newtheorem{theorem}{Theorem}
\title{Four permutation-fair dice with a three-face die}
\author{Exact construction and certificate, 30 September 2026}
\date{}
\begin{document}
\maketitle
\noindent\textbf{Status.} Verified twice by independent exact-arithmetic
programs, including direct integer subsequence counting. This is an
existence result, not an optimization of the other three dice.

\begin{theorem}
There are four permutation-fair uniform dice with face counts
\[
 \begin{split}
 K_A&=163693699141790847787594200000,\\
 K_B&=3055345920,\qquad K_C=5508,\qquad K_D=3.
 \end{split}
\]
Consequently, the minimum possible size of a prescribed die in a
four-die permutation-fair set is exactly three, using the universal
$n-1$ lower bound established in the companion monotone-tensor note.
\end{theorem}

\paragraph{The exact certificate.}
The table below specifies a run-length word. A row labelled $(i,k)$ means
$k$ consecutive occurrences of die $i$; the rows are concatenated in order.
Number all letters in this expanded word consecutively, and assign each
number to the indicated die. All labels are distinct. There are exactly
three occurrences of $D$.

For verification without expanding the very large runs, start with
$C_{\varnothing}=1$ and all other partial-pattern counts zero. Reading a
run $i^k$, simultaneously update
\[
 C_{\pi i}\ \leftarrow\ C_{\pi i}+kC_\pi
 \qquad(i\notin\pi).
\]
No pattern has a repeated letter, so it can select at most one occurrence
from the current run; this recurrence is exact. Applied to the table, it
gives for every $\pi\in S_4$
\[
 C_\pi=\frac{K_AK_BK_CK_D}{24}.
\]
The independent program \texttt{verify\_simple\_three\_face\_construction.py}
checks this identity by integer arithmetic, and separately checks all
weighted order probabilities by rational arithmetic. The machine-readable
run table is \texttt{four\_dice\_three\_face\_simple\_integer.json}.

\paragraph{How the example was found.}
At the three faces of $D$, the cumulative probabilities of the old dice
are
\[
\begin{array}{c|ccc}
 &D\text{ face 1}&D\text{ face 2}&D\text{ face 3}\\\hline
 A&1/7&71/140&17/20\\
 B&5/33&47/96&907/1056\\
 C&133/918&77/153&23/27.
\end{array}
\]
Their centered versions $y_j=2p_j-1$ satisfy
\[
 \sum_qy_j(q)=0,\qquad
 \sum_qy_i(q)y_j(q)=1\ (i\ne j),\qquad
 \sum_qy_A(q)y_B(q)y_C(q)=0.
\]
These identities give each subset of dice the probability prescribed by
a uniform ranking of being exactly the set below $D$. They alone do not
ensure the full order law. Four interval blocks
were then filled with positive rational weighted dice. A linear program
solved the remaining order constraints; rational elimination converted its
solution into an exact certificate before any assertion of fairness.

The construction uses the classical three-die fair word
\[
 \mathtt{abccba\ bcaacb\ cabbac}
\]
as a local starting point, recorded in Purcell,
\emph{Building Sets of Permutation-Fair Dice}, Example 5,
\href{https://doi.org/10.1080/00029890.2024.2345576}
{doi:10.1080/00029890.2024.2345576}.
The four rational density tilts used for $B$ are
\[
 -\frac{121}{95},\quad\frac{39391}{83895},\quad
 -\frac{319}{7020},\quad-\frac25.
\]
Die $C$ has six equal-weight atoms within each interval. The weights of
$A$ are the exact solution of the final linear system. Clearing the three
old dice's denominators gives the integer word below and leaves $D$ with
three faces.

\small
\begin{longtable}{rcr}
\text{Run}&\text{Die}&\text{Length}\\\hline
\endfirsthead
\text{Run}&\text{Die}&\text{Length}\\\hline
\endhead
1&A&1457525804001828318048582655\\
2&B&28019520\\
3&A&5891126255200470770758794782\\
4&C&133\\
5&A&312233865893543192502225816\\
6&C&133\\
7&A&777545070923506526991072450\\
8&B&60776640\\
9&A&783148640385362733842852349\\
10&B&60776640\\
11&A&777545070923506526991072450\\
12&C&133\\
13&A&917140392154551310254949284\\
14&C&133\\
15&A&777545070923506526991072450\\
16&B&93533760\\
17&A&777545070923506526991072450\\
18&C&133\\
19&A&3151103708479473819911188350\\
20&B&109912320\\
21&A&777545070923506526991072450\\
22&B&109912320\\
23&A&777545070923506526991072450\\
24&C&133\\
25&A&6207265071456708948105572064\\
26&D&1\\
27&A&706544614483005757601374890\\
28&B&212567706\\
29&A&10788535011066461264007101442\\
30&C&329\\
31&A&895698472294965317331540681\\
32&C&329\\
33&A&777545070923506526991072450\\
34&B&185624262\\
35&A&5016949959657333152877376734\\
36&B&185624262\\
37&A&5043402870558576020335777302\\
38&C&329\\
39&A&738653884932274504567638147\\
40&C&329\\
41&A&31776220877404439372527786104\\
42&B&158680818\\
43&A&777545070923506526991072450\\
44&C&329\\
45&A&777545070923506526991072450\\
46&B&145209096\\
47&A&777545070923506526991072450\\
48&B&145209096\\
49&A&777545070923506526991072450\\
50&C&329\\
51&A&777545070923506526991072450\\
52&D&1\\
53&A&3974274190980916216428391812\\
54&B&183792795\\
55&A&781729313513001301148880816\\
56&C&320\\
57&A&732178531447067349175653486\\
58&C&320\\
59&A&17110902371291397319237221726\\
60&B&186641465\\
61&A&777545070923506526991072450\\
62&C&320\\
63&A&777545070923506526991072450\\
64&B&188065800\\
65&A&777545070923506526991072450\\
66&B&188065800\\
67&A&8705230920360437285344259556\\
68&C&320\\
69&A&777545070923506526991072450\\
70&C&320\\
71&A&777545070923506526991072450\\
72&B&190914470\\
73&A&777545070923506526991072450\\
74&B&190914470\\
75&A&19376423167413782652617525454\\
76&C&320\\
77&A&777545070923506526991072450\\
78&D&1\\
79&A&770997322957834893079568682\\
80&B&57480624\\
81&A&777545070923506526991072450\\
82&C&136\\
83&A&777545070923506526991072450\\
84&C&136\\
85&A&13832117577481326638051709900\\
86&B&67060728\\
87&A&777545070923506526991072450\\
88&B&67060728\\
89&A&777545070923506526991072450\\
90&C&136\\
91&A&777545070923506526991072450\\
92&C&136\\
93&A&777545070923506526991072450\\
94&B&76640832\\
95&A&777545070923506526991072450\\
96&C&136\\
97&A&777545070923506526991072450\\
98&B&81430884\\
99&A&777545070923506526991072450\\
100&B&81430884\\
101&A&2175489261594400367097126918\\
102&C&136\\
103&A&777545070923506526991072450\\

\end{longtable}
\end{document}
